\ifdefined\gkver\else\def\gkver{student}\fi
\documentclass[\gkver]{gaokaozhenti}
\begin{document}

\chapter{2011年上海}

\begin{enumerate}
\item
%% number: 1
%% paperName: 2011年上海高考第 1 题 2 分
%% typeId: 1
%% type: 单选题
%% chapter: 电场
%% point: 电场线
%% method: 无
%% score: 2
%% degree: 850
%% duplicateId: 0
%% body:
电场线分布如图所示，电场中 a，b 两点的电场强度大小分别为已知 $E_{a}$ 和 $E_{b}$，电势分别为 $\varphi_{a}$ 和 $\varphi_{b}$，则\xzanswer{C}
\onepicture[w=4.5cm]{\includesvg{figs/01}}
\fourchoices[answer=C]
{$E_{a} > E_{b}$，$\varphi_{a} > \varphi_{b}$}
{$E_{a} > E_{b}$，$\varphi_{a} < \varphi_{b}$}
{$E_{a} < E_{b}$，$\varphi_{a} > \varphi_{b}$}
{$E_{a} < E_{b}$，$\varphi_{a} < \varphi_{b}$}
%% answer: C
%% memo:
\memoanswer{
电场线的疏密程度表示场强大小，沿电场线的方向电势逐渐降低，故 C 正确。
}

\item
%% number: 2
%% paperName: 2011年上海高考第 2 题 2 分
%% typeId: 1
%% type: 单选题
%% chapter: 原子和原子核
%% point: 原子核式结构
%% method: 无
%% score: 2
%% degree: 850
%% duplicateId: 0
%% body:
卢瑟福利用 $\alpha$ 粒子轰击金箔的实验研究原子结构，正确反映实验结果的示意图是\xzanswer{D}
\fourchoices[answer=D, ispicture=true, h=2.0cm]
{figs/02a.png}
{figs/02b.png}
{figs/02c.png}
{figs/02d.png}
%% answer: D
%% memo:
\memoanswer{
由原子核式结构，结合 $\alpha$ 粒子散射实验的现象分析：绝大部分 $\alpha$ 粒子穿过金箔后速度方向不改变，只有少数离核较近的粒子发生了较大角度的偏转（离核越近，偏转角越大），极少数正对着金原子核的 $\alpha$ 粒子被弹回，故 D 正确。
}

\item
%% number: 3
%% paperName: 2011年上海高考第 3 题 2 分
%% typeId: 1
%% type: 单选题
%% chapter: 光学
%% point: 光电效应
%% method: 无
%% score: 2
%% degree: 850
%% duplicateId: 0
%% body:
用一束紫外线照射某金属时不能产生光电效应，可能使该金属产生光电效应的措施是\xzanswer{B}
\fourchoices[answer=B]
{改用频率更小的紫外线照射}
{改用 X 射线照射}
{改用强度更大的原紫外线照射}
{延长原紫外线的照射时间}
%% answer: B
%% memo:
\memoanswer{
每种金属都有极限频率，只有入射光的频率大于金属的极限频率时才能产生光电效应，是否发生光电效应与光照强度、光照时间无关，故 B 正确。
}

\item
%% number: 4
%% paperName: 2011年上海高考第 4 题 2 分
%% typeId: 1
%% type: 单选题
%% chapter: 气体性质
%% point: 气体图象
%% method: 无
%% score: 2
%% degree: 750
%% duplicateId: 0
%% body:
如图，一定量的理想气体从状态 a 沿直线变化到状态 b，在此过程中，其压强\xzanswer{A}
\onepicture[w=4.5cm]{\includesvg{figs/04}}
\fourchoices[answer=A]
{逐渐增大}
{逐渐减小}
{始终不变}
{先增大后减小}
%% answer: A
%% memo:
\memoanswer{
过 $a$ 点和 $b$ 点分别画出等压线如图所示，从图中可以看出 $p_{a} < p_{b}$，故 A 正确。

\onepicture[h=3.8cm, align=right]{figs/04_an.jpg}
}

\item
%% number: 5
%% paperName: 2011年上海高考第 5 题 2 分
%% typeId: 1
%% type: 单选题
%% chapter: 机械振动
%% point: 单摆
%% method: 无
%% score: 2
%% degree: 650
%% duplicateId: 0
%% body:
两个相同的单摆静止于平衡位置，使摆球分别以水平初速 $v_{1}$、$v_{2}$（$v_{1} > v_{2}$）在竖直平面内做小角度摆动，它们的频率与振幅分别为 $f_{1}$、$f_{2}$ 和 $A_{1}$、$A_{2}$，则\xzanswer{C}
\fourchoices[answer=C]
{$f_{1} > f_{2}$，$A_{1} = A_{2}$}
{$f_{1} < f_{2}$，$A_{1} = A_{2}$}
{$f_{1} = f_{2}$，$A_{1} > A_{2}$}
{$f_{1} = f_{2}$，$A_{1} < A_{2}$}
%% answer: C
%% memo:
\memoanswer{
两单摆做小角度振动，由 $T = 2\pi\sqrt{\frac{L}{g}}$ 可知，两单摆摆长相同，故周期和频率都相同；由机械能守恒知，在最低点水平速率越大，向上振动得越高，振幅越大，故 C 正确。
}

\item
%% number: 6
%% paperName: 2011年上海高考第 6 题 2 分
%% typeId: 1
%% type: 单选题
%% chapter: 电路
%% point: 逻辑电路
%% method: 无
%% score: 2
%% degree: 650
%% duplicateId: 0
%% body:
右表是某逻辑电路的真值表，该电路是\xzanswer{D}

\begin{center}
\begin{tabular}{|c|c|c|}
\hline
\multicolumn{2}{|c|}{输入} & 输出 \\ \hline
0 & 0 & 1 \\ \hline
0 & 1 & 1 \\ \hline
1 & 0 & 1 \\ \hline
1 & 1 & 0 \\ \hline
\end{tabular}
\end{center}

\fourchoices[answer=D, ispicture=true, h=1.8cm]
{figs/06a.png}
{figs/06b.png}
{figs/06c.png}
{figs/06d.png}
%% answer: D
%% memo:
\memoanswer{
从真值表可以看出，此逻辑电路为“与非”门电路，故 D 正确。
}

\item
%% number: 7
%% paperName: 2011年上海高考第 7 题 2 分
%% typeId: 1
%% type: 单选题
%% chapter: 原子和原子核
%% point: 天然放射性
%% method: 无
%% score: 2
%% degree: 650
%% duplicateId: 0
%% body:
在存放放射性元素时，若把放射性元素①置于大量水中；②密封于铅盒中；③与轻核元素结合成化合物，则\xzanswer{D}
\fourchoices[answer=D]
{措施①可减缓放射性元素衰变}
{措施②可减缓放射性元素衰变}
{措施③可减缓放射性元素衰变}
{上述措施均无法减缓放射性元素衰变}
%% answer: D
%% memo:
\memoanswer{
放射性元素的衰变是从原子核内部发生的，任何物理的、化学的因素都不能改变放射性元素的衰变快慢，故 D 正确。
}

\item
%% number: 8
%% paperName: 2011年上海高考第 8 题 2 分
%% typeId: 1
%% type: 单选题
%% chapter: 气体性质
%% point: 分子动理论
%% method: 无
%% score: 2
%% degree: 650
%% duplicateId: 0
%% body:
某种气体在不同温度下的气体分子速率分布曲线如图所示，图中 $f(v)$ 表示 $v$ 处单位速率区间内的分子数百分率，所对应的温度分别为 $T_{\text{Ⅰ}}$，$T_{\text{Ⅱ}}$，$T_{\text{Ⅲ}}$，则\xzanswer{B}
\onepicture[w=4.5cm]{\includesvg{figs/08}}
\fourchoices[answer=B]
{$T_{\text{Ⅰ}} > T_{\text{Ⅱ}} > T_{\text{Ⅲ}}$}
{$T_{\text{Ⅲ}} > T_{\text{Ⅱ}} > T_{\text{Ⅰ}}$}
{$T_{\text{Ⅱ}} > T_{\text{Ⅰ}}$，$T_{\text{Ⅱ}} > T_{\text{Ⅲ}}$}
{$T_{\text{Ⅰ}} = T_{\text{Ⅱ}} = T_{\text{Ⅲ}}$}
%% answer: B
%% memo:
\memoanswer{
气体分子的平均速率随温度的升高而增大，表现出“中间多两头小”的规律，但随着温度的升高，$f(v)$ 的最大值向速率大的方向移动，故 B 正确。
}

\item
%% number: 9
%% paperName: 2011年上海高考第 9 题 3 分
%% typeId: 1
%% type: 单选题
%% chapter: 原子和原子核
%% point: 天然放射性
%% method: 无
%% score: 3
%% degree: 850
%% duplicateId: 0
%% body:
天然放射性元素放出的三种射线的穿透能力实验结果如图所示，由此可推知\xzanswer{D}
\onepicture[w=5.0cm]{\includesvg{figs/09}}
\fourchoices[answer=D]
{②来自于原子核外的电子}
{①的电离作用最强，是一种电磁波}
{③的电离作用较强，是一种电磁波}
{③的电离作用最弱，属于原子核内释放的光子}
%% answer: D
%% memo:
\memoanswer{
从图中所示三种射线的穿透能力看出，①为 $\alpha$ 射线，②为 $\beta$ 射线，③为 $\gamma$ 射线。$\beta$ 射线是从原子核变化产生的，实质是一个中子衰变成质子放出的高速电子流；$\alpha$ 射线实质上是氦核流，电离作用最强；$\gamma$ 射线为波长很短的电磁波，电离作用很弱，穿透能力最强，故 D 项正确。
}

\item
%% number: 10
%% paperName: 2011年上海高考第 10 题 3 分
%% typeId: 1
%% type: 单选题
%% chapter: 机械波
%% point: 干涉衍射
%% method: 无
%% score: 3
%% degree: 650
%% duplicateId: 0
%% body:
两波源 S$_{1}$、S$_{2}$ 在水槽中形成的波形如图所示，其中实线表示波峰，虚线表示波谷，则\xzanswer{B}
\onepicture[w=5.0cm]{\includesvg{figs/10}}
\fourchoices[answer=B]
{在两波相遇的区域中会产生干涉}
{在两波相遇的区域中不会产生干涉}
{a 点的振动始终加强}
{a 点的振动始终减弱}
%% answer: B
%% memo:
\memoanswer{
波在同种均匀介质中的波速相等，由波长、频率和波速的关系 $v = \lambda f$ 知，由于两列波的波长不同，故频率不同，所以不能产生干涉现象，故 B 正确。
}

\item
%% number: 11
%% paperName: 2011年上海高考第 11 题 3 分
%% typeId: 1
%% type: 单选题
%% chapter: 直线运动
%% point: 运动的分解
%% method: 无
%% score: 3
%% degree: 650
%% duplicateId: 0
%% body:
如图，人沿平直的河岸以速度 $v$ 行走，且通过不可伸长的绳拖船，船沿绳的方向行进，此过程中绳始终与水面平行。当绳与河岸的夹角为 $\alpha$，船的速率为\xzanswer{C}
\onepicture[w=4.5cm]{\includesvg{figs/11}}
\fourchoices[answer=C]
{$v\sin\alpha$}
{$\frac{v}{\sin\alpha}$}
{$v\cos\alpha$}
{$\frac{v}{\cos\alpha}$}
%% answer: C
%% memo:
\memoanswer{
将人的速度分解为垂直于绳的分速度 $v_{1}$ 和沿绳方向的分速度 $v_{2}$，其中沿绳方向速度大小等于船的速度，即等于 $v\cos\alpha$，故 C 正确。

\onepicture[h=3.6cm, align=right]{figs/11_an.jpg}
}

\item
%% number: 12
%% paperName: 2011年上海高考第 12 题 3 分
%% typeId: 1
%% type: 单选题
%% chapter: 电路
%% point: 动态分析
%% method: 极值
%% score: 3
%% degree: 650
%% duplicateId: 0
%% body:
如图所示电路中，闭合电键 S，当滑动变阻器的滑动触头 P 从最高端向下滑动时，\xzanswer{A}
\onepicture[w=4.5cm]{\includesvg{figs/12}}
\fourchoices[answer=A]
{电压表 V 读数先变大后变小，电流表 A 读数变大}
{电压表 V 读数先变小后变大，电流表 A 读数变小}
{电压表 V 读数先变大后变小，电流表 A 读数先变小后变大}
{电压表 V 读数先变小后变大，电流表 A 读数先变大后变小}
%% answer: A
%% memo:
\memoanswer{
滑动变阻器下部电阻 $R_{1}$ 与上部电阻 $R_{2}$ 并联。当滑片 P 下滑时，$R_{2}$ 增大，$R_{1}$ 减小，电路中的总电阻先变大后变小，电压表测路端电压，电压表示数先变大后变小。当滑片 P 下滑时，电流表先测通过 $R_{1}$ 的电流，电流表示数由零变大，$R_{1}$ 为零时，电流表测干路电流，示数最大，故 A 正确。
}

\item
%% number: 13
%% paperName: 2011年上海高考第 13 题 3 分
%% typeId: 1
%% type: 单选题
%% chapter: 电磁感应
%% point: 楞次定律推广
%% method: 无
%% score: 3
%% degree: 650
%% duplicateId: 0
%% body:
如图，均匀带正电的绝缘圆环 a 与金属圆环 b 同心共面放置，当 a 绕 O 点在其所在平面内旋转时，b 中产生顺时针方向的感应电流，且具有收缩趋势，由此可知，圆环 a\xzanswer{B}
\onepicture[w=4.0cm]{\includesvg{figs/13}}
\fourchoices[answer=B]
{顺时针加速旋转}
{顺时针减速旋转}
{逆时针加速旋转}
{逆时针减速旋转}
%% answer: B
%% memo:
\memoanswer{
b 有收缩趋势说明 a 环中电流在减小，转速减慢，使得穿过 b 环的磁通量减小，所以 b 环中感应电流的方向与 a 环方向相同，故 B 正确。
}

\item
%% number: 14
%% paperName: 2011年上海高考第 14 题 3 分
%% typeId: 1
%% type: 单选题
%% chapter: 电场
%% point: 电场中的图象
%% method: 无
%% score: 3
%% degree: 650
%% duplicateId: 0
%% body:
两个等量异种点电荷位于 $x$ 轴上，相对原点对称分布，正确描述电势 $\varphi$ 随位置 $x$ 变化规律的是图\xzanswer{A}
\fourchoices[answer=A, ispicture=true, h=2.6cm]
{figs/14a.png}
{figs/14b.png}
{figs/14c.png}
{figs/14d.png}
%% answer: A
%% memo:
\memoanswer{
画出等量异种点电荷周围电场线如图所示，其中虚线所示的中垂线为等势线。由于沿电场线方向电势逐渐降低，且连线的中点与无限远处等势，故 A 正确。

\onepicture[h=3.8cm, align=right]{figs/14_an.jpg}
}

\item
%% number: 15
%% paperName: 2011年上海高考第 15 题 3 分
%% typeId: 1
%% type: 单选题
%% chapter: 功和能
%% point: 功率
%% method: 无
%% score: 3
%% degree: 650
%% duplicateId: 0
%% body:
如图，一长为 $L$ 的轻杆一端固定在光滑铰链上，另一端固定一质量为 $m$ 的小球。一水平向右的拉力作用于杆的中点，使杆以角速度 $\omega$ 匀速转动，当杆与水平方向成 $60^{\circ}$ 时，拉力的功率为\xzanswer{C}
\onepicture[w=4.0cm]{\includesvg{figs/15}}
\fourchoices[answer=C]
{$mgL\omega$}
{$\frac{\sqrt{3}}{2}mgL\omega$}
{$\frac{1}{2}mgL\omega$}
{$\frac{\sqrt{3}}{6}mgL\omega$}
%% answer: C
%% memo:
\memoanswer{
轻杆中点速度大小为 $v = \frac{1}{2}\omega L$ ①；由有固定转动轴的物体的平衡条件得 $mgL\cos 60^{\circ} = F\cdot\frac{1}{2}L\sin 60^{\circ}$ ②；拉力的功率 $P = Fv\cos 30^{\circ}$ ③，由①②③得 $P = \frac{1}{2}mgL\omega$，故 C 正确。
}

\item
%% number: 16
%% paperName: 2011年上海高考第 16 题 3 分
%% typeId: 1
%% type: 单选题
%% chapter: 运动和力
%% point: 连接体
%% method: 无
%% score: 3
%% degree: 450
%% duplicateId: 0
%% body:
如图，在水平面上的箱子内，带异种电荷的小球 a、b 用绝缘细线分别系于上、下两边，处于静止状态。地面受到的压力为 $N$，球 b 所受细线的拉力为 $F$。剪断连接球 b 的细线后，在球 b 上升过程中地面受到的压力\xzanswer{D}
\onepicture[w=4.0cm]{\includesvg{figs/16}}
\fourchoices[answer=D]
{小于 $N$}
{等于 $N$}
{等于 $N + F$}
{大于 $N + F$}
%% answer: D
%% memo:
\memoanswer{
剪断细线前，系统处于静止状态，压力 $N$ 等于系统的总重力，即 $N = Mg + m_{a}g + m_{b}g$；剪断连接球 b 的细线后，b 球向上做加速运动，系统处于超重状态，压力大于总重力，故 AB 错误。

再选 a 和箱为研究对象，箱子受到的支持力 $N' = Mg + m_{a}g + F_{\text{库}}$，其中 $F_{\text{库}}$ 为 b 球对 a 球的库仑力。因为平衡时 $F_{\text{库}} = m_{b}g + F$，而球 b 加速上升过程中，随着两球距离的减小，$F_{\text{库}}' > F_{\text{库}} = m_{b}g + F$，故 $N' > Mg + m_{a}g + m_{b}g + F = N + F$，故 D 正确。
}

\item
%% number: 17
%% paperName: 2011年上海高考第 17 题 4 分
%% typeId: 2
%% type: 多选题
%% chapter: 光学
%% point: 波粒二象性
%% method: 无
%% score: 4
%% degree: 850
%% duplicateId: 0
%% body:
用极微弱的可见光做双缝干涉实验，随着时间的增加，在屏上先后出现如图~\ref{2011上海17a}、\ref{2011上海17b}、\ref{2011上海17c} 所示的图像，则\xzanswer{ABD}
\threepicture[labelA=2011上海17a, labelB=2011上海17b, labelC=2011上海17c, hA=2.2cm, hB=2.2cm, hC=2.2cm]
{figs/17a.png}
{figs/17b.png}
{figs/17c.png}
\fourchoices[answer=ABD]
{图像（a）表明光具有粒子性}
{图像（c）表明光具有波动性}
{用紫外光观察不到类似的图像}
{实验表明光是一种概率波}
%% answer: ABD
%% memo:
\memoanswer{
照射时间短时，在光屏上显示出一些亮点，说明光子表现出粒子性；当照射时间比较长时，单个光子到达屏的位置虽然是不确定的，但大量光子到达光屏上的位置服从统计规律，表明大量光子表现出波动性，说明光是一种概率波，故 ABD 正确。
}

\item
%% number: 18
%% paperName: 2011年上海高考第 18 题 4 分
%% typeId: 2
%% type: 多选题
%% chapter: 磁场
%% point: 安培力
%% method: 无
%% score: 4
%% degree: 650
%% duplicateId: 0
%% body:
如图，质量为 $m$、长为 $L$ 的直导线用两绝缘细线悬挂于 O、O$'$，并处于匀强磁场中。当导线中通以沿 $x$ 正方向的电流 $I$，且导线保持静止时，悬线与竖直方向夹角为 $\theta$。则磁感应强度方向和大小可能为\xzanswer{BC}
\onepicture[w=4.0cm]{\includesvg{figs/18}}
\fourchoices[answer=BC]
{$z$ 正向，$\frac{mg}{IL}\tan\theta$}
{$y$ 正向，$\frac{mg}{IL}$}
{$z$ 负向，$\frac{mg}{IL}\tan\theta$}
{沿悬线向上，$\frac{mg}{IL}\sin\theta$}
%% answer: BC
%% memo:
\memoanswer{
由左手定则先判断各种情况下安培力的方向，看是否平衡，若平衡再由平衡条件求解 $B$ 的大小。若磁感应强度 $B$ 沿 $z$ 正向，安培力方向向左，不可能平衡，故 A 错误；若 $B$ 沿 $y$ 正向，安培力方向竖直向上，可以平衡，由平衡条件得 $BIL = mg$，可得 $B = \frac{mg}{IL}$，故 B 正确。

若 $B$ 沿 $z$ 负向，安培力方向水平向右，受力分析如图所示，则有 $mg\tan\theta = BIL$，可得 $B = \frac{mg}{IL}\tan\theta$，故 C 正确；若 $B$ 沿悬线向上则安培力方向斜向左下，不能平衡，故 D 错误。

\onepicture[h=3.8cm, align=right]{figs/18_an.jpg}
}

\item
%% number: 19
%% paperName: 2011年上海高考第 19 题 4 分
%% typeId: 2
%% type: 多选题
%% chapter: 运动和力
%% point: 变加速直线运动
%% method: 无
%% score: 4
%% degree: 650
%% duplicateId: 0
%% body:
受水平外力 $F$ 作用的物体，在粗糙水平面上作直线运动，其 $v-t$ 图线如图所示，则\xzanswer{CD}
\onepicture[w=4.5cm]{\includesvg{figs/19}}
\fourchoices[answer=CD]
{在 $0 \sim t_{1}$ 秒内，外力 $F$ 大小不断增大}
{在 $t_{1}$ 时刻，外力 $F$ 为零}
{在 $t_{1} \sim t_{2}$ 秒内，外力 $F$ 大小可能不断减小}
{在 $t_{1} \sim t_{2}$ 秒内，外力 $F$ 大小可能先减小后增大}
%% answer: CD
%% memo:
\memoanswer{
从 $v-t$ 图象可以看出，加速度先为正，大小减小；后为负，大小增大。$t_{1}$ 时刻加速度为 0。由牛顿第二定律可知，$t_{1}$ 时刻合外力为零，故 B 错误；在 $0 \sim t_{1}$ 秒内，$F - f = ma$，$a$ 减小，$F$ 减小，故 A 错误；在 $t_{1} \sim t_{2}$ 秒内，当外力 $F$ 与阻力 $f$ 方向相反时，$f - F = ma$，$a$ 增大，外力 $F$ 减小；当外力 $F$ 与阻力 $f$ 方向相同时，$f + F = ma$，$a$ 增大，外力 $F$ 增大，故 CD 正确。
}

\item
%% number: 20
%% paperName: 2011年上海高考第 20 题 4 分
%% typeId: 2
%% type: 多选题
%% chapter: 电磁感应
%% point: 楞次定律推广
%% method: 无
%% score: 4
%% degree: 450
%% duplicateId: 0
%% body:
如图，磁场垂直于纸面，磁感应强度在竖直方向均匀分布，水平方向非均匀分布。一铜制圆环用丝线悬挂于 O 点，将圆环拉至位置 a 后无初速释放，在圆环从 a 摆向 b 的过程中\xzanswer{AD}
\onepicture[w=4.5cm]{\includesvg{figs/20}}
\fourchoices[answer=AD]
{感应电流方向先逆时针后顺时针再逆时针}
{感应电流方向一直是逆时针}
{安培力方向始终与速度方向相反}
{安培力方向始终沿水平方向}
%% answer: AD
%% memo:
\memoanswer{
本题磁场的特点是靠近中心线磁场较强，而两侧磁场较弱。圆环从 a 摆向 b 的过程中，磁感线先是向里增多，感应电流的方向是逆时针；然后是向里减少，向外增多，感应电流的方向是顺时针；再后是向外减小，感应电流的方向是逆时针。由于上、下半圆受到安培力大小相等，所以只需分析左、右半圆受到的安培力即可。安培力总是阻碍圆环由左向右运动，故方向水平向左，故 AD 正确。
}

\item
%% number: 21
%% paperName: 2011年上海高考第 21 题 4 分
%% typeId: 3
%% type: 填空题
%% chapter: 光学
%% point: 光的衍射
%% method: 无
%% score: 4
%% degree: 850
%% duplicateId: 0
%% body:
如图，当用激光照射直径小于激光束的不透明圆盘时，在圆盘后屏上的阴影中心出现了一个亮斑。这是光的\tkanswer[1.2cm]{衍射}（填“干涉”、“衍射”或“直线传播”）现象，这一实验支持了光的\tkanswer[1.5cm]{波动说}（填“波动说”、“微粒说”或“光子说”）。
\onepicture[w=6.0cm, align=right]{figs/21.png}
%% answer: 衍射，波动说
%% memo:
\memoanswer{
光在传播过程中，一般呈现直线传播的特征，但是当障碍物的尺度与光的波长可比拟时，将出现干涉或衍射现象。题意给出在屏上的阴影中心出现了一个亮斑，这是衍射产生的条纹。干涉、衍射现象支持了光的波动说。本题的正确答案为：衍射；波动说。
}

\item
%% number: 22
%% paperName: 2011年上海高考第 22 题 4 分
%% typeId: 3
%% type: 填空题
%% chapter: 运动和力
%% point: 动量守恒定律
%% method: 无
%% score: 4
%% degree: 850
%% duplicateId: 0
%% body:
光滑水平面上两小球 a、b 用不可伸长的松弛细绳相连。开始时 a 球静止，b 球以一定速度运动直至绳被拉紧，然后两球一起运动，在此过程中两球的总动量\tkanswer[1.1cm]{守恒}（填“守恒”或“不守恒”）；机械能\tkanswer[1.1cm]{不守恒}（填“守恒”或“不守恒”）。
%% answer: 守恒，不守恒
%% memo:
\memoanswer{
两球相互作用的过程满足动量守恒定律的条件——系统受到的合外力为零；但不满足机械能守恒定律的条件，因为在绳被拉紧的短暂过程中，机械能有损失转化为内能（这类似于完全非弹性碰撞）。
}

\item
%% number: 22
%% paperName: 2011年上海高考第 22 题 4 分
%% typeId: 3
%% type: 填空题
%% chapter: 曲线运动
%% point: 人造地球卫星
%% method: 无
%% score: 4
%% degree: 650
%% duplicateId: 0
%% body:
人造地球卫星在运行过程中由于受到微小的阻力，轨道半径将缓慢减小。在此运动过程中，卫星所受万有引力大小将\tkanswer[1.1cm]{增大}（填“减小”或“增大”）；其动能将\tkanswer[1.1cm]{增大}（填“减小”或“增大”）。
%% answer: 增大，增大
%% memo:
\memoanswer{
由万有引力定律 $F = G\frac{Mm}{r^{2}}$ 知，随着卫星距地心的距离逐渐减小，引力将增大；由于轨道半径逐渐减小，故每一状态都可以看成是圆周运动，由 $G\frac{Mm}{r^{2}} = m\frac{v^{2}}{r}$ 得 $E_{k} = \frac{1}{2}mv^{2} = \frac{GM}{2r}$，所以随着 $r$ 的减小，动能将增大。
}

\item
%% number: 23
%% paperName: 2011年上海高考第 23 题 4 分
%% typeId: 3
%% type: 填空题
%% chapter: 电场
%% point: 带电粒子的圆周运动
%% method: 无
%% score: 4
%% degree: 650
%% duplicateId: 0
%% body:
如图，在竖直向下，场强为 $E$ 的匀强电场中，长为 $l$ 的绝缘轻杆可绕固定轴在竖直面内无摩擦转动，两个小球 A、B 固定于杆的两端，A、B 的质量分别为 $m_{1}$ 和 $m_{2}$（$m_{1} < m_{2}$），A 带负电，电量为 $q_{1}$，B 带正电，电量为 $q_{2}$。杆从静止开始由水平位置转到竖直位置，在此过程中电场力做功为\tkanswer[2.2cm]{$\frac{(q_{1}+q_{2})El}{2}$}，在竖直位置处两球的总动能为\tkanswer[3.8cm]{$\frac{(q_{1}+q_{2})El+(m_{2}-m_{1})gl}{2}$}。
\onepicture[align=right]{\includesvg{figs/23}}
%% answer: $\frac{(q_{1}+q_{2})El}{2}$，$\frac{(q_{1}+q_{2})El+(m_{2}-m_{1})gl}{2}$
%% memo:
\memoanswer{
A 球受到的电场力方向向上、B 球受到的电场力方向向下，电场力对两球均做正功，电场力做的总功

$W = q_{1}E\cdot\frac{l}{2} + q_{2}E\cdot\frac{l}{2} = \frac{(q_{1}+q_{2})El}{2}$

两球从静止开始运动，故两球的总动能等于合外力对系统做的功

$\Delta E_{k} = \frac{(q_{1}+q_{2})El}{2} + m_{2}g\frac{l}{2} - m_{1}g\frac{l}{2} = \frac{\left[(q_{1}+q_{2})E+(m_{2}-m_{1})g\right]l}{2}$。
}

\item
%% number: 24
%% paperName: 2011年上海高考第 24 题 4 分
%% typeId: 3
%% type: 填空题
%% chapter: 机械波
%% point: 波的叠加
%% method: 无
%% score: 4
%% degree: 650
%% duplicateId: 0
%% body:
两列简谐波沿 $x$ 轴相向而行，波速均为 $v = 0.4\Ums$，两波源分别位于 A、B 处，$t = 0$ 时的波形如图所示。当 $t = 2.5\Us$ 时，M 点的位移为\tkanswer[0.8cm]{2}$\Ucm$，N 点的位移为\tkanswer[0.8cm]{0}$\Ucm$。
\onepicture[align=right]{\includesvg{figs/24}}
%% answer: 2，0
%% memo:
\memoanswer{
从图中可以看出，两列波的振幅均为 $2\Ucm$，由 $v = \frac{\lambda}{T}$ 可知，A 波的周期 $T_{A} = 0.5\Us$，B 波的周期 $T_{B} = 1\Us$。则经过 $2.5\Us$ 的时间，A 波向右传播了 5 个周期，B 波向左传播 2.5 个周期，作出波形图知，两列波在 N 点均处于平衡位置，A 波在 M 点处的质点处于平衡位置，B 波位于波峰位置，故 M 点位移为 $2\Ucm$，N 点位移为 0。
}

\item
%% number: 25
%% paperName: 2011年上海高考第 25 题 4 分
%% typeId: 3
%% type: 填空题
%% chapter: 曲线运动
%% point: 平抛运动
%% method: 无
%% score: 4
%% degree: 450
%% duplicateId: 0
%% body:
以初速为 $v_{0}$，射程为 $s$ 的平抛运动轨迹制成一光滑轨道。一物体由静止开始从轨道顶端滑下，当其到达轨道底部时，物体的速率为\tkanswer[1.6cm]{$\frac{gs}{v_{0}}$}，其水平方向的速度大小为\tkanswer[3.4cm]{$\frac{v_{0}}{\sqrt{1+\left(\frac{v_{0}^{2}}{gs}\right)^{2}}}$}。
%% answer: $\frac{gs}{v_{0}}$，$\frac{v_{0}}{\sqrt{1+\left(\frac{v_{0}^{2}}{gs}\right)^{2}}}$
%% memo:
\memoanswer{
由平抛运动的规律可得轨道高度为 $h = \frac{1}{2}g\left(\frac{s}{v_{0}}\right)^{2} = \frac{gs^{2}}{2v_{0}^{2}}$。物体从轨道顶端由静止滑到底部时，由机械能守恒定律得 $mgh = \frac{1}{2}mv^{2}$，得 $v = \frac{gs}{v_{0}}$。

设速度方向与水平方向的夹角为 $\theta$，由平抛运动的规律可得 $\tan\theta = \frac{v_{y}}{v_{x}} = \frac{\sqrt{2gh}}{v_{0}} = \frac{gs}{v_{0}^{2}}$；而物块落到轨道底部时水平分速度 $v_{x} = v\cos\theta$，解得 $v_{x} = \frac{v_{0}}{\sqrt{1+\left(\frac{v_{0}^{2}}{gs}\right)^{2}}}$。
}

\item
%% number: 26
%% paperName: 2011年上海高考第 26 题 5 分
%% typeId: 4
%% type: 实验题
%% chapter: 直线运动
%% point: DIS测瞬时速度
%% method: 无
%% score: 5
%% degree: 550
%% duplicateId: 0
%% body:
如图，为测量作匀加速直线运动小车的加速度，将宽度均为 $b$ 的挡光片 A、B 固定在小车上，测得二者间距为 $d$。
\begin{enumerate}
	\item
	当小车匀加速经过光电门时，测得两挡光片先后经过的时间 $\Delta t_{1}$ 和 $\Delta t_{2}$，则小车加速度 $a = $\tkanswer[2.6cm]{$\frac{b^{2}}{2d}\left(\frac{1}{(\Delta t_{2})^{2}} - \frac{1}{(\Delta t_{1})^{2}}\right)$}。
	\item
	（多选题）为减小实验误差，可采取的方法是（\quad）

	（A）增大两挡光片宽度 $b$　　（B）减小两挡光片宽度 $b$

	（C）增大两挡光片间距 $d$　　（D）减小两挡光片间距 $d$
\end{enumerate}
\onepicture[align=right]{\includesvg{figs/26}}
%% answer: （1）$\frac{b^{2}}{2d}\left(\frac{1}{(\Delta t_{2})^{2}} - \frac{1}{(\Delta t_{1})^{2}}\right)$（2）BC
\jdanswer{
\begin{enumerate}
	\item
	$\frac{b^{2}}{2d}\left(\frac{1}{(\Delta t_{2})^{2}} - \frac{1}{(\Delta t_{1})^{2}}\right)$

	\item
	BC
\end{enumerate}
}
%% memo:
\memoanswer{
（1）物块通过两光电门的速度分别为 $v_{a} = \frac{b}{\Delta t_{1}}$，$v_{b} = \frac{b}{\Delta t_{2}}$，由 $v_{b}^{2} - v_{a}^{2} = 2ad$ 得 $a = \frac{b^{2}}{2d}\left[\frac{1}{(\Delta t_{2})^{2}} - \frac{1}{(\Delta t_{1})^{2}}\right]$。

（2）本实验测速度的原理是用挡光片通过光电门时的平均速度代替瞬时速度，所以挡光片通过光电门的时间越短，即宽度越小，误差越小；另外两挡光片间的距离越大，误差越小，故 BC 正确。
}

\item
%% number: 27
%% paperName: 2011年上海高考第 27 题 5 分
%% typeId: 4
%% type: 实验题
%% chapter: 气体性质
%% point: 用单分子油膜估测分子大小
%% method: 无
%% score: 5
%% degree: 650
%% duplicateId: 0
%% body:
在“用单分子油膜估测分子大小”实验中：
\begin{enumerate}
	\item
	某同学操作步骤如下：
	\begin{steps}
		\item
		取一定量的无水酒精和油酸，制成一定浓度的油酸酒精溶液；
		\item
		在量筒中滴入一滴该溶液，测出它的体积；
		\item
		在蒸发皿内盛一定量的水，再滴入一滴油酸酒精溶液，待其散开稳定；
		\item
		在蒸发皿上覆盖透明玻璃，描出油膜形状，用透明方格纸测量油膜的面积。
	\end{steps}
	改正其中的错误：\tkanswer[4.6cm]{②在量筒中滴入 $N$ 滴溶液；③在水面上先撒上痱子粉}。
	\item
	若油酸酒精溶液体积浓度为 0.10\%，一滴溶液的体积为 $4.8 \times 10^{-3}\UmL$，其形成的油膜面积为 $40\Ucmq$，则估测出油酸分子的直径为\tkanswer[2.0cm]{$1.2 \times 10^{-9}$}$\Um$。
\end{enumerate}
%% answer: （1）②在量筒中滴入 $N$ 滴溶液；③在水面上先撒上痱子粉（2）$1.2 \times 10^{-9}$
\jdanswer{
\begin{enumerate}
	\item
	②在量筒中滴入 $N$ 滴溶液；③在水面上先撒上痱子粉

	\item
	$1.2 \times 10^{-9}$
\end{enumerate}
}
%% memo:
\memoanswer{
（1）步骤②不正确，因为滴入一滴溶液很难测得该滴溶液的精确体积，有读数误差。在实验中通常是滴入 $N$ 滴该溶液，测定 $N$ 滴溶液的体积时读数的误差将会减小。在操作步骤③前应先在水面上撒上痱子粉，然后将该溶液滴在痱子粉上使得容易区分溶液与水。

（2）一滴溶液含油酸的体积 $V = 4.8 \times 10^{-3}\UmL \times 0.10\% = 4.8 \times 10^{-6}\UmL$，由 $V = Sd$ 可得，油酸的直径 $d = \frac{V}{S} = 1.2 \times 10^{-9}\Um$。
}

\item
%% number: 28
%% paperName: 2011年上海高考第 28 题 5 分
%% typeId: 1
%% type: 单选题
%% chapter: 电磁感应
%% point: 感应电动势与磁通量变化快慢的关系
%% method: 无
%% score: 5
%% degree: 650
%% duplicateId: 0
%% body:
在“研究回路中感应电动势大小与磁通量变化快慢的关系”实验（见图~\ref{2011上海28a}）中，得到 $E - \frac{1}{\Delta t}$ 图线如图~\ref{2011上海28b} 所示。
\twopicture[labelA=2011上海28a, labelB=2011上海28b, hA=3.4cm, hB=4.4cm]
{figs/28a.png}
{figs/28b.png}
（1）（多选题）在实验中需保持不变的是\xzanswer{AD}
\fourchoices[answer=AD]
{挡光片的宽度}
{小车的释放位置}
{导轨倾斜的角度}
{光电门的位置}
（2）线圈匝数增加一倍后重做该实验，在图~\ref{2011上海28b} 中画出实验图线。
\drawpicanswer{\onepicture[h=4.4cm]{figs/28b.png}}{\onepicture[h=4.4cm]{\includesvg{figs/28_an}}}
%% answer: AD
%% memo:
\memoanswer{
（1）由法拉第电磁感应定律知 $E = n\frac{\Delta\Phi}{\Delta t}$，从图象中可以看出 $E \propto \frac{1}{\Delta t}$，说明实验中控制 $\Delta\Phi$ 不变，所以应保持光电门的位置和挡光片的宽度不变，故 AD 正确。

（2）$E - \frac{1}{\Delta t}$ 的斜率等于 $n\Delta\Phi$，在 $\Delta\Phi$ 不变时，若 $n$ 增加一倍，则斜率增加一倍，图线见答图。
}

\item
%% number: 29
%% paperName: 2011年上海高考第 29 题 9 分
%% typeId: 4
%% type: 实验题
%% chapter: 电路
%% point: 电路实验
%% method: 无
%% score: 9
%% degree: 450
%% duplicateId: 0
%% body:
实际电流表有内阻，可等效为理想电流表与电阻的串联。测量实际电流表 $G_{1}$ 内阻 $r_{1}$ 的电路如图所示。供选择的仪器如下：①待测电流表 $G_{1}$（$0 \sim 5\UmA$，内阻约 $300\UO$），②电流表 $G_{2}$（$0 \sim 10\UmA$，内阻约 $100\UO$），③定值电阻 $R_{1}$（$300\UO$），④定值电阻 $R_{2}$（$10\UO$），⑤滑动变阻器 $R_{3}$（$0 \sim 1000\UO$），⑥滑动变阻器 $R_{4}$（$0 \sim 20\UO$），⑦干电池（$1.5\UV$），⑧电键 S 及导线若干。
\onepicture[h=4.0cm, align=right]{\includesvg{figs/29a}}
\begin{enumerate}
	\item
	定值电阻应选\tkanswer[0.8cm]{③}，滑动变阻器应选\tkanswer[0.8cm]{⑥}。（在空格内填写序号）
	\item
	用连线连接实物图。
	\drawpicanswer{\onepicture[h=4.6cm]{\includesvg{figs/29b}}}{\onepicture[h=4.6cm]{\includesvg{figs/29_an}}}
	\item
	补全实验步骤：
	\begin{steps}
		\item
		按电路图连接电路，\tkanswer[3.8cm]{将滑动触头移至最左端}；
		\item
		闭合电键 S，移动滑动触头至某一位置，记录 $G_{1}$，$G_{2}$ 的读数 $I_{1}$，$I_{2}$；
		\item
		\tkanswer[5.2cm]{多次移动滑动触头，记录相应的 $G_{1}$，$G_{2}$ 读数 $I_{1}$，$I_{2}$}；
		\item
		以 $I_{2}$ 为纵坐标，$I_{1}$ 为横坐标，作出相应图线，如图所示。
		\onepicture[h=3.6cm, align=right]{\includesvg{figs/29c}}
	\end{steps}
	\item
	根据 $I_{2} - I_{1}$ 图线的斜率 $k$ 及定值电阻，写出待测电流表内阻的表达式\tkanswer[2.4cm]{$r_{1} = (k - 1)R_{1}$}。
\end{enumerate}
%% answer: （1）③，⑥（2）见图（3）①将滑动触头移至最左端；③多次移动滑动触头，记录相应的 $G_{1}$，$G_{2}$ 读数 $I_{1}$，$I_{2}$（4）$r_{1} = (k - 1)R_{1}$
\jdanswer{
\begin{enumerate}
	\item
	③，⑥

	\item
	见图

	\item
	①将滑动触头移至最左端；③多次移动滑动触头，记录相应的 $G_{1}$，$G_{2}$ 读数 $I_{1}$，$I_{2}$

	\item
	$r_{1} = (k - 1)R_{1}$
\end{enumerate}
}
%% memo:
\memoanswer{
（1）由实验原理图，待测电阻 $r_{1} = \frac{I_{2} - I_{1}}{I_{1}}R$，因为两电表的量程之差与待测电表的量程相当，所以与 $G_{1}$ 并联的定值电阻的阻值与 $r_{1}$ 差不多，约为 $300\UO$，故定值电阻选 $R_{1}$；滑动变阻器用的是分压式接法，应选用阻值较小的变阻器，选 $R_{4}$。

（4）由 $r_{1} = \frac{I_{2} - I_{1}}{I_{1}}R_{1} = \left(\frac{I_{2}}{I_{1}} - 1\right)R_{1}$ 得 $r_{1} = (k - 1)R_{1}$。
}

\item
%% number: 30
%% paperName: 2011年上海高考第 30 题 10 分
%% typeId: 6
%% type: 计算题
%% chapter: 气体性质
%% point: 气体连接体
%% method: 无
%% score: 10
%% degree: 650
%% duplicateId: 0
%% body:
如图，绝热气缸 A 与导热气缸 B 均固定于地面，由刚性杆连接的绝热活塞与两气缸间均无摩擦。两气缸内装有处于平衡状态的理想气体，开始时体积均为 $V_{0}$、温度均为 $T_{0}$。缓慢加热 A 中气体，停止加热达到稳定后，A 中气体压强为原来的 1.2 倍。设环境温度始终保持不变，求气缸 A 中气体的体积 $V_{A}$ 和温度 $T_{A}$。
\onepicture[align=right]{\includesvg{figs/30}}
%% answer: $V_{A} = \frac{7}{6}V_{0}$，$T_{A} = 1.4T_{0}$
\jdanswer{
$V_{A} = \frac{7}{6}V_{0}$，$T_{A} = 1.4T_{0}$
}
%% memo:
\memoanswer{
设初态压强为 $p_{0}$，膨胀后 A、B 压强相等，有

$p_{B} = 1.2p_{0}$

B 中气体始末状态温度相等，由玻意耳定律得

$p_{0}V_{0} = 1.2p_{0}(2V_{0} - V_{A})$

解得 $V_{A} = \frac{7}{6}V_{0}$。

A 部分气体满足

$\frac{p_{0}V_{0}}{T_{0}} = \frac{1.2p_{0}V_{A}}{T_{A}}$

解得 $T_{A} = 1.4T_{0}$。
}

\item
%% number: 31
%% paperName: 2011年上海高考第 31 题 12 分
%% typeId: 6
%% type: 计算题
%% chapter: 运动和力
%% point: 牛顿定律的应用
%% method: 无
%% score: 12
%% degree: 650
%% duplicateId: 0
%% body:
如图，质量 $m = 2\Ukg$ 的物体静止于水平地面的 A 处，A、B 间距 $L = 20\Um$。用大小为 $30\UN$，沿水平方向的外力拉此物体，经 $t_{0} = 2\Us$ 拉至 B 处。（已知 $\cos 37^{\circ} = 0.8$，$\sin 37^{\circ} = 0.6$。取 $g = 10\Umsq$）
\begin{enumerate}
	\item
	求物体与地面间的动摩擦因数 $\mu$；
	\item
	用大小为 $30\UN$，与水平方向成 $37^{\circ}$ 的力斜向上拉此物体，使物体从 A 处由静止开始运动并能到达 B 处，求该力作用的最短时间 $t$。
\end{enumerate}
\onepicture[align=right]{\includesvg{figs/31}}
%% answer: （1）$\mu = 0.5$（2）$t = 1.03\Us$
\jdanswer{
\begin{enumerate}
	\item
	$\mu = 0.5$

	\item
	$t = 1.03\Us$
\end{enumerate}
}
%% memo:
\memoanswer{
（1）物体做匀加速运动，由运动学公式得

$L = \frac{1}{2}at_{0}^{2}$

解得 $a = \frac{2L}{t_{0}^{2}} = \frac{2 \times 20}{2^{2}}\Umsq = 10\Umsq$。

由牛顿第二定律

$F - f = ma$

解得 $f = (30 - 2 \times 10)\UN = 10\UN$。

故 $\mu = \frac{f}{mg} = \frac{10}{2 \times 10} = 0.5$。

（2）设 $F$ 作用的最短时间为 $t$，物体先以大小为 $a$ 的加速度匀加速 $t$ 秒，撤去外力后，以大小为 $a'$ 的加速度匀减速 $t'$ 秒到达 B 处，速度恰为 0，由牛顿第二定律

$F\cos 37^{\circ} - \mu(mg - F\sin 37^{\circ}) = ma$

解得 $a = \frac{F(\cos 37^{\circ} + \mu\sin 37^{\circ})}{m} - \mu g = \left(\frac{30 \times (0.8 + 0.5 \times 0.6)}{2} - 0.5 \times 10\right)\Umsq = 11.5\Umsq$。

$a' = \frac{f}{m} = \mu g = 5\Umsq$

由于匀加速阶段的末速度即为匀减速阶段的初速度，因此有

$at = a't'$

解得 $t' = \frac{a}{a'}t = \frac{11.5}{5}t = 2.3t$。

又 $L = \frac{1}{2}at^{2} + \frac{1}{2}a't'^{2}$

解得 $t = \sqrt{\frac{2L}{a + 2.3^{2}a'}} = \sqrt{\frac{2 \times 20}{11.5 + 2.3^{2} \times 5}}\Us = 1.03\Us$。

（2）另解：设力 $F$ 作用的最短时间为 $t$，相应的位移为 $s$，物体到达 B 处速度恰为 0，由动能定理

$\left[F\cos 37^{\circ} - \mu(mg - F\sin 37^{\circ})\right]s - \mu mg(L - s) = 0$

解得 $s = \frac{\mu mgL}{F(\cos 37^{\circ} + \mu\sin 37^{\circ})} = \frac{0.5 \times 2 \times 10 \times 20}{30 \times (0.8 + 0.5 \times 0.6)}\Um = 6.06\Um$。

由牛顿第二定律 $F\cos 37^{\circ} - \mu(mg - F\sin 37^{\circ}) = ma$

解得 $a = \frac{F(\cos 37^{\circ} + \mu\sin 37^{\circ})}{m} - \mu g = 11.5\Umsq$

又 $s = \frac{1}{2}at^{2}$

解得 $t = \sqrt{\frac{2s}{a}} = \sqrt{\frac{2 \times 6.06}{11.5}}\Us = 1.03\Us$。
}

\item
%% number: 32
%% paperName: 2011年上海高考第 32 题 14 分
%% typeId: 5
%% type: 辨析题
%% chapter: 电磁感应
%% point: 电磁感应中的变加速
%% method: 无
%% score: 14
%% degree: 550
%% duplicateId: 0
%% body:
电阻可忽略的光滑平行金属导轨长 $s = 1.15\Um$，两导轨间距 $L = 0.75\Um$，导轨倾角为 $30^{\circ}$，导轨上端 ab 接一阻值 $R = 1.5\UO$ 的电阻，磁感应强度 $B = 0.8\UT$ 的匀强磁场垂直轨道平面向上。阻值 $r = 0.5\UO$，质量 $m = 0.2\Ukg$ 的金属棒与轨道垂直且接触良好，从轨道上端 ab 处由静止开始下滑至底端，在此过程中金属棒产生的焦耳热 $Q_{r} = 0.1\UJ$。（取 $g = 10\Umsq$）求：
\begin{enumerate}
	\item
	金属棒在此过程中克服安培力的功 $W_{\text{安}}$；
	\item
	金属棒下滑速度 $v = 2\Ums$ 时的加速度 $a$；
	\item
	为求金属棒下滑的最大速度 $v_{m}$，有同学解答如下：由动能定理 $W_{\text{重}} - W_{\text{安}} = \frac{1}{2}mv_{m}^{2}$，$\ldots\ldots$。由此所得结果是否正确？若正确，说明理由并完成本小题；若不正确，给出正确的解答。
\end{enumerate}
\onepicture[align=right]{\includesvg{figs/32}}
%% answer: （1）$W_{\text{安}} = 0.4\UJ$（2）$a = 3.2\Umsq$（3）此解法正确，$v_{m} = 2.74\Ums$
\jdanswer{
\begin{enumerate}
	\item
	$W_{\text{安}} = 0.4\UJ$

	\item
	$a = 3.2\Umsq$

	\item
	此解法正确，$v_{m} = 2.74\Ums$
\end{enumerate}
}
%% memo:
\memoanswer{
（1）下滑过程中安培力的功即为在电阻上产生的焦耳热，由于 $R = 3r$，因此

$Q_{R} = 3Q_{r} = 0.3\UJ$

故 $W_{\text{安}} = Q = Q_{R} + Q_{r} = 0.4\UJ$。

（2）金属棒下滑时受重力和安培力

$F_{\text{安}} = BIL = \frac{B^{2}L^{2}v}{R + r}$

由牛顿第二定律 $mg\sin 30^{\circ} - \frac{B^{2}L^{2}v}{R + r} = ma$

解得 $a = g\sin 30^{\circ} - \frac{B^{2}L^{2}v}{m(R + r)} = \left(10 \times 0.5 - \frac{0.8^{2} \times 0.75^{2} \times 2}{0.2 \times (1.5 + 0.5)}\right)\Umsq = 3.2\Umsq$。

（3）此解法正确。

金属棒下滑时受重力和安培力作用，其运动满足

$mg\sin 30^{\circ} - \frac{B^{2}L^{2}v}{R + r} = ma$

上式表明，加速度随速度增加而减小，棒做加速度减小的加速运动。无论最终是否达到匀速，当棒到达斜面底端时速度一定为最大。由动能定理可以得到棒的末速度，因此上述解法正确。

$mgs\sin 30^{\circ} - Q = \frac{1}{2}mv_{m}^{2}$

解得 $v_{m} = \sqrt{2gs\sin 30^{\circ} - \frac{2Q}{m}} = \sqrt{2 \times 10 \times 1.15 \times 0.5 - \frac{2 \times 0.4}{0.2}}\Ums = 2.74\Ums$。
}

\item
%% number: 33
%% paperName: 2011年上海高考第 33 题 14 分
%% typeId: 6
%% type: 计算题
%% chapter: 功和能
%% point: 能量守恒定律
%% method: 无
%% score: 14
%% degree: 450
%% duplicateId: 0
%% body:
如图~\ref{2011上海33a}，磁铁 A、B 的同名磁极相对放置，置于水平气垫导轨上。A 固定于导轨左端，B 的质量 $m = 0.5\Ukg$，可在导轨上无摩擦滑动。将 B 在 A 附近某一位置由静止释放，由于能量守恒，可通过测量 B 在不同位置处的速度，得到 B 的势能随位置 $x$ 的变化规律，见图~\ref{2011上海33c} 中曲线 Ⅰ。若将导轨右端抬高，使其与水平面成一定角度（如图~\ref{2011上海33b} 所示），则 B 的总势能曲线如图~\ref{2011上海33c} 中 Ⅱ 所示，将 B 在 $x = 20\Ucm$ 处由静止释放，求：（解答时必须写出必要的推断说明。取 $g = 9.8\Umsq$）
\begin{enumerate}
	\item
	B 在运动过程中动能最大的位置；
	\item
	运动过程中 B 的最大速度和最大位移。
	\item
	图~\ref{2011上海33c} 中直线 Ⅲ 为曲线 Ⅱ 的渐近线，求导轨的倾角。
	\item
	若 A、B 异名磁极相对放置，导轨的倾角不变，在图~\ref{2011上海33c} 上画出 B 的总势能随 $x$ 的变化曲线。
\end{enumerate}
\threepicture[labelA=2011上海33a, labelB=2011上海33b, labelC=2011上海33c, hA=2.4cm, hB=2.4cm, hC=4.6cm, align=right]
{figs/33a.png}
{figs/33b.png}
{figs/33c.png}
%% answer: （1）$x = 6.1\Ucm$（2）$v_{m} = 1.31\Ums$，$\Delta x = 18.0\Ucm$（3）$\theta = 59.7^{\circ}$（4）如图
\jdanswer{
\begin{enumerate}
	\item
	$x = 6.1\Ucm$

	\item
	$v_{m} = 1.31\Ums$，$\Delta x = 18.0\Ucm$

	\item
	$\theta = 59.7^{\circ}$

	\item
	如图
\end{enumerate}
}
%% memo:
\memoanswer{
（1）由能量守恒知，势能最小处动能最大。由图线 Ⅱ 得 $x = 6.1\Ucm$（在 $5.9 \sim 6.3\Ucm$ 间均视为正确）。

（2）由图读得释放处势能 $E_{p} = 0.90\UJ$，此即 B 的总能量。

由于运动中总能量守恒，因此在势能最小处动能最大，由图像得最小势能为 $0.47\UJ$，则最大动能为

$E_{km} = 0.9 - 0.47 = 0.43\UJ$（$E_{km}$ 在 $0.42 \sim 0.44\UJ$ 间均视为正确）

最大速度为

$v_{m} = \sqrt{\frac{2E_{km}}{m}} = \sqrt{\frac{2 \times 0.43}{0.5}}\Ums = 1.31\Ums$（$v_{m}$ 在 $1.29 \sim 1.33\Ums$ 间均视为正确）

$x = 20.0\Ucm$ 处的总能量为 $0.90\UJ$，最大位移由 $E = 0.90\UJ$ 的水平直线与曲线 Ⅱ 的左侧交点确定，由图中读出交点位置为 $x = 2.0\Ucm$，因此，最大位移

$\Delta x = 20.0 - 2.0 = 18.0\Ucm$（$\Delta x$ 在 $17.9 \sim 18.1\Ucm$ 间均视为正确）

（3）渐近线 Ⅲ 表示 B 的重力势能随位置变化关系，即

$E_{pg} = mgx\sin\theta = kx$

解得 $\sin\theta = \frac{k}{mg}$

由图读出直线斜率为

$k = \frac{0.85 - 0.30}{20.0 - 7.0} = 4.23 \times 10^{-2}\UJ[a]/\Ucm[a]$

$\theta = \sin^{-1}\left(\frac{k \times 10^{2}}{mg}\right) = \sin^{-1}\left(\frac{4.23}{0.5 \times 9.8}\right) = 59.7^{\circ}$（$\theta$ 在 $59^{\circ} \sim 61^{\circ}$ 间均视为正确）

（4）若异名磁极相对放置，A、B 间相互作用势能为负值，总势能如图。

\onepicture[h=4.6cm, align=right]{\includesvg{figs/33_an}}
}

\end{enumerate}

\end{document}
